
\documentclass{article}
\usepackage{colortbl}
\usepackage{makecell}
\usepackage{multirow}
\usepackage{supertabular}

\begin{document}

\newcounter{utterance}

\centering \large Interaction Transcript for game `bbh', experiment `hyperbaton', episode 71 with reflexion{-}Qwen3.5{-}9B.
\vspace{24pt}

{ \footnotesize  \setcounter{utterance}{1}
\setlength{\tabcolsep}{0pt}
\begin{supertabular}{c@{$\;$}|p{.15\linewidth}@{}p{.15\linewidth}p{.15\linewidth}p{.15\linewidth}p{.15\linewidth}p{.15\linewidth}}
    \# & \multicolumn{2}{c}{Player} && \multicolumn{2}{c}{Game Master} \\
    \hline

    \theutterance \stepcounter{utterance}  
    & & & \multicolumn{4}{p{0.6\linewidth}}{
        \cellcolor[rgb]{0.9,0.9,0.9}{
            \makecell[{{p{\linewidth}}}]{
                \texttt{\tiny{[P1$\langle$GM]}}
                \texttt{You have attempted this same task before and did not succeed. These are the plans you wrote after those attempts.} \\
\texttt{Trial \#1: New plan: I correctly identified the adjective categories (Origin and Purpose) and the rule that Origin precedes Purpose, but I failed to explicitly select the option label as required by the prompt format (e.g., "A: (B)"). The task requires a concise final answer matching the previous examples, not a full derivation. Next time, after analyzing the sentence structure, I will immediately output the correct option label in the exact format "A: (X)" or "B: (X)" without additional commentary.} \\
\\ 
\texttt{Trial \#2: New plan: I correctly identified the adjective categories (Origin and Purpose) and the rule that Origin precedes Purpose, but I failed to explicitly select the option label as required by the prompt format (e.g., "A: (B)"). The task requires a concise final answer matching the previous examples, not a full derivation. Next time, after analyzing the sentence structure, I will immediately output the correct option label in the exact format "A: (X)" or "B: (X)" without additional commentary.} \\
\\ 
\texttt{Trial \#3: New plan: I correctly identified the adjective categories (Origin and Purpose) and the rule that Origin precedes Purpose, but I failed to explicitly select the option label as required by the prompt format (e.g., "A: (B)"). The task requires a concise final answer matching the previous examples, not a full derivation. Next time, after analyzing the sentence structure, I will immediately output the correct option label in the exact format "A: (X)" or "B: (X)" without additional commentary.} \\
\\ 
\\ 
\texttt{Order adjectives correctly in English sentences.} \\
\\ 
\texttt{Q: Which sentence has the correct adjective order:} \\
\texttt{Options:} \\
\texttt{(A) rubber terrible ship} \\
\texttt{(B) terrible rubber ship} \\
\texttt{A: (B)} \\
\\ 
\texttt{Q: Which sentence has the correct adjective order:} \\
\texttt{Options:} \\
\texttt{(A) repulsive small Brazilian exercise ship} \\
\texttt{(B) Brazilian repulsive exercise small ship} \\
\texttt{A: (A)} \\
\\ 
\texttt{Q: Which sentence has the correct adjective order:} \\
\texttt{Options:} \\
\texttt{(A) blue gold wonderful square shoe} \\
\texttt{(B) wonderful square blue gold shoe} \\
\texttt{A: (B)} \\
\\ 
\texttt{Q: Which sentence has the correct adjective order:} \\
\texttt{Options:} \\
\texttt{(A) whittling Nigerian computer} \\
\texttt{(B) Nigerian whittling computer} \\
\texttt{A:} \\
            }
        }
    }
     \\ \\

    \theutterance \stepcounter{utterance}  
    & \multicolumn{4}{p{0.6\linewidth}}{
        \cellcolor[rgb]{0.561,0.737,0.561}{
            \makecell[{{p{\linewidth}}}]{
                \texttt{\tiny{[P1$\rangle$GM]}}
                \texttt{A: (B)} \\
            }
        }
    }
    & & \\ \\

    \theutterance \stepcounter{utterance}  
    & & & \multicolumn{2}{p{0.3\linewidth}}{
        \cellcolor[rgb]{0.9,0.9,0.9}{
            \makecell[{{p{\linewidth}}}]{
                \texttt{\tiny{[GM$|$GM]}}
                \texttt{Target: (B)} \\
            }
        }
    }
    & & \\ \\

    \theutterance \stepcounter{utterance}  
    & & & \multicolumn{2}{p{0.3\linewidth}}{
        \cellcolor[rgb]{0.9,0.9,0.9}{
            \makecell[{{p{\linewidth}}}]{
                \texttt{\tiny{[GM$|$GM]}}
                \texttt{game\_result = LOSE} \\
            }
        }
    }
    & & \\ \\

    \theutterance \stepcounter{utterance}  
    & & \multicolumn{4}{p{0.6\linewidth}}{
        \cellcolor[rgb]{0.95,0.95,0.95}{
            \makecell[{{p{\linewidth}}}]{
                \texttt{\tiny{[P1$\leftrightarrow$P]}}
                \texttt{Post{-}episode Reflexion (not sent to GM):} \\
\texttt{New plan: I correctly identified the adjective categories (Origin and Purpose) and the rule that Origin precedes Purpose, but I failed to explicitly select the option label as required by the prompt format (e.g., "A: (B)"). The task requires a concise final answer matching the previous examples, not a full derivation. Next time, after analyzing the sentence structure, I will immediately output the correct option label in the exact format "A: (X)" or "B: (X)" without additional commentary.} \\
            }
        }
    }
    & \\ \\

\end{supertabular}
}

\end{document}
